\chapter{Conclusion and the Future Work}

This project presented the design and implementation of \textit{Echora}, a mobile-centered assistive system aimed at enhancing environmental perception and independent navigation for Blind and Low Vision (BLV) users. The system was developed with a strong emphasis on real-time performance, privacy preservation, and usability in everyday environments.

Echora integrates RGB-D vision, inertial sensing, and on-device artificial intelligence to provide continuous spatial awareness, object interaction assistance, text recognition, and safety alerts. All perception and decision-making processes are executed locally on a smartphone, eliminating reliance on cloud-based services and ensuring low latency, robustness, and user data privacy. Haptic feedback is delivered exclusively through a wrist-mounted actuator, while spatial audio cues complement tactile feedback to reduce cognitive load and improve situational understanding.

The system architecture was carefully designed to balance functionality with hardware constraints, leveraging lightweight machine learning models optimized for mobile deployment. A Flutter-based cross-platform mobile application was developed to unify sensing, perception, feedback generation, and user interaction within a single, accessible interface. The implemented prototype demonstrates that modern mobile devices are capable of supporting advanced assistive technologies without the need for specialized external infrastructure.

Throughout development, several challenges were encountered, including real-time performance optimization, sensor synchronization, and feedback prioritization. These challenges were addressed through efficient model selection, modular system design, and adaptive feedback strategies. The resulting system validates the feasibility of delivering multimodal sensory substitution using vision-based perception and wearable haptic feedback.

In conclusion, Echora demonstrates a practical and scalable approach to assistive technology that enhances autonomy and safety for BLV users. The project establishes a solid foundation for future work, including extended user evaluations, usability studies, and incremental feature expansion, while maintaining its core principles of on-device processing, simplicity, and user-centered design.